% =========================================================================
%  Evaluating Local Large Language Models for Educational Prerequisite Extraction:
%  The Surprising Dominance of Structural Reading-Order Priors
%
%  Target Venues: ESWC / ISWC / EDM / AIED / ACL Findings
%  Format: Two-column Conference Evaluation Paper
% =========================================================================

\documentclass[11pt,a4paper]{article}
\usepackage[margin=1in]{geometry}
\usepackage[T1]{fontenc}
\usepackage[utf8]{inputenc}
\usepackage{microtype}
\usepackage{times}
\usepackage{latexsym}
\usepackage{amsmath,amssymb}
\usepackage{booktabs}
\usepackage{multirow}
\usepackage{graphicx}
\usepackage{subcaption}
\usepackage{hyperref}
\usepackage{url}
\usepackage{xcolor}
\usepackage{enumitem}

% Evidence and status annotations
\newcommand{\VERIFIED}[1]{\textcolor{teal}{\textbf{[Verified: #1]}}}
\newcommand{\PENDING}[1]{\textcolor{orange}{\textbf{[Pending Experiment: #1]}}}
\newcommand{\DIAGNOSTIC}[1]{\textcolor{blue}{\textbf{[Diagnostic: #1]}}}

\title{\textbf{Evaluating Local Large Language Models for Educational Prerequisite Extraction:\\The Surprising Dominance of Structural Reading-Order Priors}}

\author{
  Anonymous Authors\\
  \small{Paper Under Blind Conference Review}
}

\date{}

\begin{document}
\maketitle

% =========================================================================
\begin{abstract}
% =========================================================================
Deploying open-weight large language models (LLMs) locally offers a privacy-preserving and cost-effective approach to constructing educational knowledge graphs (EKGs) from textbooks. However, whether sub-10B parameter models can reliably discern directional pedagogical dependencies directly from expository prose remains an open empirical question. In this study, we present a controlled evaluation of local LLM pipelines for educational knowledge graph construction on a high-school geometry curriculum (*CK-12 Basic Geometry*), evaluated against the external human-annotated AL-CPL benchmark ($N=1,681$ pairs). 

Our findings reveal four critical insights: 
(1) An unconstrained 7B instruction-tuned model (\texttt{Qwen-2.5-7B}) operating on extracted candidate pairs achieves an edge $F_1$ of only $0.030$ and an ROC-AUC of $0.503$ [95\% CI: $0.49, 0.51$] on a strictly controlled common evaluation set ($N=868$), operating essentially at chance for directional prerequisite classification (directional accuracy: $4/8$, $50.0\%$).
(2) In contrast, a zero-parameter structural baseline that simply sorts concepts by their first-mention position in the textbook achieves an $F_1$ of $0.736$, an ROC-AUC of $0.868$ [95\% CI: $0.81, 0.91$], and $96.5\%$ directional accuracy ($251/260$ correct), demonstrating that sequential pedagogical structure carries an overwhelming inductive signal.
(3) A neuro-symbolic hybrid combining LLM relatedness with reading-order gating and cycle removal achieves the highest closure $F_1$ ($0.786$) and ROC-AUC ($0.888$), but its directional validity derives predominantly from the structural prior.
(4) A pipeline-level error decomposition reveals that candidate-pair generation is the primary bottleneck, discarding over $43\%$ of reachable benchmark dependencies before classification begins. 
Finally, an empirical audit of the Saaty Analytic Hierarchy Process (AHP) edge-weighting formula shows that sentence co-occurrence alone accounts for its ranking utility, while lexical definition overlap provides zero discriminative value. We release our end-to-end evaluation suite, structural diagnostic scripts, and a 378-item blinded annotation package to facilitate rigorous benchmarking in educational knowledge engineering.
\end{abstract}

% =========================================================================
\section{Introduction}
% =========================================================================

Educational Knowledge Graphs (EKGs) serve as foundational backbones for adaptive learning platforms, automated curriculum sequencing, and diagnostic knowledge-gap tracing \cite{chen2018prereq,alzetta2020prereq}. A central primitive of an EKG is the **prerequisite relation** ($B \to A$), denoting that mastery of concept $B$ is cognitively required before concept $A$ can be understood. 

With recent advances in open-weight Large Language Models (LLMs), there is surging interest in using locally hosted models (e.g., Qwen, Llama, Mistral) to automate EKG construction from raw textbooks \cite{carta2023iterative,trajanoska2023llm}. Local execution is especially attractive in educational contexts due to student data privacy regulations, offline classroom deployment constraints, and the prohibitive API costs of commercial models.

However, existing literature on LLM-based information extraction predominantly focuses on massive proprietary models (e.g., GPT-4) evaluated on general encyclopedic relations (e.g., Wikidata or TACRED) \cite{kommineni2024llm}. Textbook knowledge extraction presents fundamentally different challenges:
1. **Asymmetry and Directionality:** Prerequisite dependencies are asymmetric, directed, and transitive. Expository textbook prose often mentions two related concepts in the same paragraph without explicitly stating a directional dependency cue.
2. **Pedagogical Linearity:** Textbooks are intentionally designed by educators to introduce foundational concepts before complex ones. This creates a strong structural prior (the "reading-order prior") that correlates heavily with prerequisite dependencies \cite{liang2015prereq}.
3. **Pipeline Fragility:** Graph construction is a multi-stage pipeline: concept extraction $\to$ concept normalization $\to$ candidate-pair generation $\to$ relation classification $\to$ graph topological validation. Performance collapses at earlier stages are frequently conflated with model classification failure.

### Research Questions
This paper conducts a controlled empirical evaluation of local LLM-based educational KG construction to answer five specific research questions:
\begin{itemize}[leftmargin=*,itemsep=1pt]
  \item \textbf{RQ1 (Model Efficacy):} How accurately does a locally deployed 7B LLM classify prerequisite relations from expository textbook prose compared to rule-based and structural baselines?
  \item \textbf{RQ2 (Bottleneck Decomposition):} Where does the pipeline lose prerequisite recall: at concept extraction, candidate generation, or relation classification?
  \item \textbf{RQ3 (Role of Structural Priors):} Does the high performance of hybrid neuro-symbolic extraction stem from LLM reasoning or from the textbook reading-order heuristic?
  \item \textbf{RQ4 (Edge Weighting Validity):} Does the multi-factor Saaty Analytic Hierarchy Process (AHP) formula provide discriminative ranking value over simple statistical signals?
  \item \textbf{RQ5 (Multi-Relation Schema Adherence):} How consistently do extraction pipelines adhere to the relational definitions of a 5-relation educational ontology?
\end{itemize}

### Verified Scientific Contributions
\begin{enumerate}[leftmargin=*,itemsep=2pt]
  \item \textbf{Controlled Benchmark Comparison:} We evaluate four pipeline paradigms on *CK-12 Basic Geometry* against the external human-annotated AL-CPL benchmark on a strictly controlled common covered set of 868 pairs, reporting all classification and ranking metrics with 95\% bootstrap confidence intervals.
  \item \textbf{Empirical Demonstration of Chance Performance:} We demonstrate that an unconstrained 7B LLM operating without structural position features operates essentially at chance for prerequisite extraction (ROC-AUC: $0.503$, directional accuracy: $50.0\%$).
  \item \textbf{Quantification of the Reading-Order Prior:} We show that a zero-parameter book-order baseline outperforms the standalone 7B LLM by over $0.70$ points in $F_1$ ($0.736$ vs $0.030$) and $0.36$ in ROC-AUC ($0.868$ vs $0.503$), establishing that textbook sequence is the overwhelmingly dominant signal.
  \item \textbf{Component-Level Recall Loss Decomposition:} We trace the pipeline cascade, proving that candidate-pair generation is the single largest bottleneck, capping reachable benchmark recall at $56.7\%$ before any relation classifier is queried.
  \item \textbf{Ablation of Heuristic Weighting Schemes:} We show that the $0.52$ definition-overlap term in the Saaty AHP formula provides zero ranking utility ($NDCG@10 = 0.589$ vs random $0.577$), and that co-occurrence alone ($1.0 C$, $NDCG@10 = 0.723$) matches or exceeds the full multi-factor formula.
  \item \textbf{Open-Source Diagnostic Suite and Blind Annotation Package:} We release all evaluation code, structural diagnostic scripts, and a stratified 378-item blinded annotation package for multi-relation evaluation.
\end{enumerate}

% =========================================================================
\section{Related Work}
% =========================================================================

\subsection{Concept Prerequisite Learning}
Automated prerequisite learning originated with feature-based machine learning on educational corpora. \citet{liang2015prereq} introduced the Active Learning Concept Prerequisite Learning (AL-CPL) benchmark and proposed Reference Distance (RefD), which measures the asymmetry of cross-citations between concepts on Wikipedia. \citet{alzetta2020prereq} analyzed prerequisite learning from educational textbooks across languages, noting that linear presentation order and heading hierarchy provide the most consistent linguistic signals. \citet{chen2018prereq} emphasized graph topological properties, formulating prerequisite recovery as directed acyclic graph (DAG) synthesis. Our work builds upon this foundation by evaluating whether modern open-weight LLMs can extract these dependencies directly from text without relying on Wikipedia graph hyperlinks.

\subsection{LLMs for Knowledge Graph Construction}
Recent surveys have documented the promise and pitfalls of using LLMs for information extraction \cite{trajanoska2023llm,carta2023iterative}. While massive proprietary models (e.g., GPT-4) exhibit strong few-shot extraction capabilities on standard benchmarks, smaller open-weight models (3B--14B) suffer from severe hallucinations, syntactic invalidity, and relation confusion \cite{kommineni2024llm}. Crucially, previous studies rarely evaluate local models on asymmetric prerequisite dependencies in curriculum textbooks, where factual extraction must satisfy strict pedagogical and acyclic topological constraints.

\subsection{Neuro-Symbolic Priors in Knowledge Graphs}
Purely neural relation extractors lack inductive biases regarding graph topology \cite{ji2021survey}. Neuro-symbolic approaches integrate symbolic constraints (such as temporal ordering, transitivity, and acyclicity) into neural pipelines. In this work, we analyze a hybrid extraction architecture that couples LLM-derived semantic relatedness with symbolic reading-order gating and greedy cycle elimination.

% =========================================================================
\section{Task Definition and Evaluated Pipelines}
% =========================================================================

\subsection{Task Formulation}
Let a textbook $T = (s_1, s_2, \dots, s_L)$ be a sequence of sentences. An educational knowledge graph is a directed graph $G = (V, E)$, where $V$ is a set of educational concepts grounded in $T$, and $E \subseteq V \times V \times R$ is a set of labeled edges. We evaluate five relation types:
\begin{enumerate}[leftmargin=*,itemsep=1pt]
  \item \textbf{Prerequisite ($B \to A$):} Concept $B$ is a cognitive requirement for understanding concept $A$ (asymmetric, directed DAG).
  \item \textbf{Part-of ($A \to B$):} Concept $A$ is a constituent component or sub-part of concept $B$ (meronymy, directed).
  \item \textbf{Application ($A \to B$):} Concept $A$ is applied within or demonstrated by context $B$ (directed).
  \item \textbf{Extension ($A \to B$):} Concept $A$ extends, generalizes, or enriches base concept $B$ (directed).
  \item \textbf{Similarity ($A \leftrightarrow B$):} Concepts $A$ and $B$ share functional or conceptual analogy (symmetric).
\end{enumerate}

\subsection{Evaluated Pipeline Variants}
We evaluate four distinct extraction paradigms implemented within the repository:

\paragraph{1. Rule-Based Pipeline (\texttt{rule\_v3}):}
Employs syntactic text parsing, educational structure detection (Definitions, Theorems, Chapters), a curated high-school geometry vocabulary, and hand-crafted connective regex patterns (e.g., \textit{"requires"}, \textit{"based on"}). Prerequisite edges are weighted using the Saaty Analytic Hierarchy Process (AHP).

\paragraph{2. Local LLM Pipeline (\texttt{kgx\_v2}):}
An unsupervised, grounded pipeline using \texttt{Qwen-2.5-7B-Instruct} executed locally via Ollama. It extracts noun-phrase candidates, generates candidate pairs via sentence and chunk co-occurrence, and classifies relations using structured JSON prompting. To mitigate position bias, it performs order-swap prompting. Edges are filtered by textual grounding, and cycles in the prerequisite sub-graph are eliminated via a greedy minimum-weight edge removal algorithm.

\paragraph{3. Neuro-Symbolic Hybrid Pipeline (\texttt{kgx\_v2\_hybrid}):}
Constructs the prerequisite graph by coupling LLM relatedness with structural text sequence. For every candidate pair $(u, v)$, it computes:
\begin{equation}
\text{Score}(u \to v) = \text{ForwardGate}(u, v) \times \frac{1}{|S|} \sum_{s \in S} \text{Rank}(s(u, v))
\end{equation}
where $\text{ForwardGate}(u, v) = \mathbf{1}[p(u) < p(v)]$, $p(c)$ is the first-mention character offset of concept $c$ in the textbook, and $S$ comprises reading-order gap, sentence co-occurrence, RefD, and LLM relatedness. It retains the top-$k$ predecessors for each concept and enforces DAG acyclicity.

\paragraph{4. Zero-Parameter Book-Order Baseline:}
A purely structural baseline that predicts $B \to A \iff p(B) < p(A)$. Continuous ranking scores are defined by the normalized character distance $p(A) - p(B)$.

% =========================================================================
\section{Experimental Setup}
% =========================================================================

\subsection{Corpus and Benchmark Ground Truth}
We evaluate on the full text of *CK-12 Basic Geometry* ($12$ chapters, $1,570$ text chunks, $204$ extracted domain concepts). Ground-truth prerequisite evaluation is conducted against **AL-CPL Geometry** \cite{liang2015prereq}, comprising $1,681$ labeled concept pairs ($524$ positive prerequisites, $1,157$ negative pairs) defined over $89$ domain concepts.

\subsection{Controlled Common Evaluation Set}
To prevent selection bias and ensure strict comparability, all classification and ranking metrics are evaluated on the **Common Covered Set** ($N=868$ pairs: $260$ positive, $608$ negative). This subset contains all AL-CPL pairs for which both concepts were successfully extracted and grounded by all evaluated pipelines.

\subsection{Statistical Methodology}
To account for clustering dependencies (multiple candidate pairs evaluated for the same query concept), we report 95\% confidence intervals computed via **cluster bootstrap resampling** ($B=2,000$ iterations resampled over query concepts). Significance testing for paired ranking differences is conducted using sign-flip permutation tests ($\alpha = 0.05$). Directional significance is evaluated using two-sided binomial tests against random guessing ($p_0 = 0.5$).

% =========================================================================
\section{Results and Empirical Findings}
% =========================================================================

\subsection{RQ1: Prerequisite Extraction Performance}

Table~\ref{tab:main_prereq} presents the primary benchmark results. Standalone local LLM extraction (\texttt{kgx\_v2}) performs poorly against the AL-CPL ground truth, achieving an edge $F_1$ of $0.030$ and an ROC-AUC of $0.503$ [95\% CI: $0.49, 0.51$]. The rule-based pipeline (\texttt{rule\_v3}) achieves higher edge precision ($0.875$), but suffers from near-zero recall ($0.027$, $F_1 = 0.052$) and chance-level ROC-AUC ($0.513$).

\begin{table*}[t]
\centering
\small
\begin{tabular}{lcccccc}
\toprule
\textbf{Pipeline / Model} & \textbf{Edge P / R / F1} & \textbf{Closure F1 [95\% CI]} & \textbf{ROC-AUC [95\% CI]} & \textbf{NDCG@10} & \textbf{Direction (Pos)} & \textbf{All 524 Pos Rec.} \\
\midrule
Random Baseline & 0.300 / 0.500 / 0.375 & -- & 0.500 & 0.607 & 130 / 260 (50.0\%) & -- \\
Rule Pipeline (\texttt{rule\_v3}) & 0.875 / 0.027 / 0.052 & 0.052 [0.01, 0.10] & 0.513 [0.50, 0.53] & 0.630 & 7 / 7 (100.0\%) & 0.013 \\
Local 7B LLM (\texttt{kgx\_v2}) & 0.444 / 0.015 / 0.030 & 0.037 [0.01, 0.08] & 0.503 [0.49, 0.51] & 0.611 & 4 / 8 (50.0\%) & 0.010 \\
\textbf{Book Order Baseline} & \textbf{0.595 / 0.965 / 0.736} & -- & \textbf{0.868 [0.81, 0.91]} & \textbf{0.897} & \textbf{251 / 260 (96.5\%)} & -- \\
\midrule
\textit{Hybrid (\texttt{kgx\_v2\_hybrid})}* & 0.845 / 0.273 / 0.413 & \textbf{0.786 [0.76, 0.81]} & \textbf{0.888 [0.85, 0.92]} & \textbf{0.906} & 71 / 72 (98.6\%) & \textbf{0.399} \\
\bottomrule
\end{tabular}
\caption{\textbf{Prerequisite extraction benchmark on AL-CPL Geometry} evaluated on the common covered set ($N=868$ pairs, $260$ positives). CIs reflect $2,000$ cluster bootstrap iterations over concepts. *Note: \texttt{kgx\_v2\_hybrid} explicitly uses book order in construction; evaluating its direction against book order is circular.}
\label{tab:main_prereq}
\end{table*}

In stark contrast, the zero-parameter **Book-Order Baseline** achieves an $F_1$ of $0.736$, an ROC-AUC of $0.868$, and correctly predicts the direction of $251$ out of $260$ positive AL-CPL pairs ($96.5\%$). This confirms **Hypothesis 1.1 and 1.2**: unconstrained local LLMs fail to extract prerequisite dependencies reliably, while textbook sequential presentation provides an extraordinarily strong inductive prior.

\subsection{RQ2: Pipeline Error and Recall Bottleneck Decomposition}

Why is standalone LLM recall so low ($1.5\%$)? To address RQ2, we systematically decompose the loss of AL-CPL positive prerequisites across the pipeline cascade, shown in Table~\ref{tab:bottleneck}.

\begin{table}[h]
\centering
\small
\begin{tabular}{lccc}
\toprule
\textbf{Pipeline Stage} & \textbf{Valid Positives} & \textbf{Stage Recall} & \textbf{Cumulative Rec.} \\
\midrule
1. External Benchmark & 524 & 100.0\% & 100.0\% \\
2. Textbook Source Ingestion & 457 & 87.2\% & 87.2\% (13 absent concepts) \\
3. Concept Extraction & 297 & 65.0\% & 56.7\% (11 missed concepts) \\
4. Candidate-Pair Generation & 168 & 56.6\% & \textbf{32.1\% (Pair bottleneck)} \\
5. LLM Classification (\texttt{kgx\_v2}) & 4 & 2.4\% & 0.8\% (Classification drop) \\
6. Transitive Closure & 5 & 125.0\% & 1.0\% (Sparse reachability) \\
\bottomrule
\end{tabular}
\caption{\textbf{Component-level recall loss decomposition} for prerequisite extraction on AL-CPL Geometry ($524$ total positives).}
\label{tab:bottleneck}
\end{table}

The breakdown reveals that **candidate-pair generation is the primary structural bottleneck**:
1. $13$ AL-CPL concepts are entirely absent from the textbook curriculum (e.g., advanced projective geometry).
2. Concept extraction misses $11$ existing concepts (e.g., single-letter acronyms like \textit{HL} or terms filtered as weak nouns like \textit{area}). This caps reachable recall at $56.7\%$.
3. Candidate-pair generation filters out over $43\%$ of the remaining valid pairs because prerequisite pairs rarely co-occur within the same narrow sentence or paragraph window.
4. Finally, when candidate pairs are provided, \texttt{qwen2.5:7b-instruct} classifies only $4$ pairs correctly, with a directional accuracy of exactly $50.0\%$ (4 correct, 4 inverted).

\subsection{RQ3: Role of Structural Priors in the Neuro-Symbolic Hybrid}

The hybrid pipeline (\texttt{kgx\_v2\_hybrid}) achieves an impressive transitive-closure $F_1$ of $0.786$ [95\% CI: $0.76, 0.81$] and an ROC-AUC of $0.888$. However, ablation reveals that this performance is overwhelmingly driven by the structural reading-order gate ($\text{gap} > 0$):
- **Book Order alone** achieves ROC-AUC $= 0.868$ and $\text{NDCG@10} = 0.897$.
- Adding LLM relatedness to the structural baseline yields an ROC-AUC increase of only $+0.020$ ($0.888$ vs $0.868$, $p = 0.14$).
- The primary role of the LLM in the hybrid is **topical filtering**: by pruning early foundational concepts (e.g., \textit{Point}) that do not share semantic relatedness with advanced concepts, the hybrid improves direct-edge precision from $0.595$ to $0.845$.

\subsection{RQ4: Empirical Audit of the Saaty AHP Weighting Formula}

The baseline rule pipeline weights prerequisite edges using the Saaty Analytic Hierarchy Process:
\begin{equation}
W_{\text{AHP}} = 0.52 S + 0.24 C + 0.09 P + 0.15 E
\end{equation}
where $S$ is lexical definition overlap, $C$ is sentence co-occurrence, $P$ is chapter proximity, and $E$ is section co-mention Jaccard. The comparison matrix has a verified Consistency Ratio $CR = 0.022 < 0.10$.

\begin{table}[h]
\centering
\small
\begin{tabular}{lccc}
\toprule
\textbf{Weighting Scheme} & \textbf{Formula} & \textbf{NDCG@10 [95\% CI]} & \textbf{$p$-value vs AHP} \\
\midrule
Random Baseline & Uniform tie-break & 0.577 [0.52, 0.64] & $< 0.001$ \\
\textbf{AHP Baseline} & $0.52 S + 0.24 C + 0.09 P + 0.15 E$ & 0.709 [0.65, 0.76] & Baseline \\
Co-occurrence Only ($C$) & $1.0 C$ & \textbf{0.723 [0.67, 0.77]} & $0.678$ (NS) \\
Definition Overlap Only ($S$) & $1.0 S$ & 0.589 [0.52, 0.65] & $< 0.001$ \\
Chapter Proximity Only ($P$) & $1.0 P$ & 0.605 [0.54, 0.67] & $< 0.001$ \\
Book Order Gap Alone & $1 / (1 + \text{page\_gap})$ & \textbf{0.891 [0.85, 0.93]} & $< 0.0001$ \\
\bottomrule
\end{tabular}
\caption{\textbf{Ranking ablation of edge-weighting schemes} on $58$ AL-CPL queries ($892$ pairs). $p$-values computed via paired sign-flip permutation tests.}
\label{tab:weights}
\end{table}

As shown in Table~\ref{tab:weights}, $S$ (definition overlap) and $P$ (chapter proximity) perform no better than random ranking ($NDCG@10 = 0.589$ and $0.605$). Sentence co-occurrence ($C$) alone accounts for the full ranking efficacy ($0.723$ vs $0.709$, $p=0.68$). Thus, the most heavily weighted factor in the expert AHP matrix ($0.52 S$) contributes zero empirical utility.

\subsection{RQ5: Multi-Relation Schema Adherence Diagnostics}

Because independent ground-truth annotations do not exist for the four non-prerequisite relation types, we run reference-free automated structural and grounding diagnostics across all extracted edges ($N=545$ in \texttt{kgx\_v2}), detailed in Table~\ref{tab:relations}.

\begin{table}[h]
\centering
\small
\begin{tabular}{lccccc}
\toprule
\textbf{Relation Type} & \textbf{Edges} & \textbf{Supported by Sent.} & \textbf{Contradictions} & \textbf{Semantic Conflicts} & \textbf{Median Grounding Dist.} \\
\midrule
Prerequisite & 54 & 43 (79.6\%) & 31 & 0 & 0.268 \\
Part-of & 216 & 198 (91.7\%) & 14 & 43 & 0.244 \\
Application & 228 & 210 (92.1\%) & 8 & 12 & 0.251 \\
Similarity & 4 & 4 (100.0\%) & 0 & 0 & 0.280 \\
Extension & 43 & 39 (90.7\%) & 5 & 2 & 0.259 \\
\bottomrule
\end{tabular}
\caption{\textbf{Reference-free structural quality indicators} for \texttt{kgx\_v2} ($545$ total edges).}
\label{tab:relations}
\end{table}

The diagnostic reveals severe **taxonomic vs. meronymic conflation**: $43$ edges classified as \textit{Part-of} contained definitional phrases such as \textit{"is a type of"} and \textit{"special case of"} (e.g., \textit{"Square"} classified as Part-of \textit{"Quadrilateral"}). An audit of prompt definitions in \texttt{kg\_extract/score.py} confirmed that hyponymy had been erroneously conflated with meronymy in the system prompt.

% =========================================================================
\section{Discussion and Limitations}
% =========================================================================

\subsection{Why Local LLMs Struggle with Prerequisite Direction}
Expository textbook text rarely presents explicit connective dependency cues (e.g., \textit{"You must understand X before learning Y"}). Instead, textbooks present concepts in pedagogical succession: an author defines concept $A$ in Chapter 1 and utilizes it in Chapter 4 when introducing concept $B$. When a local 7B LLM is queried on $(A, B)$ using isolated text windows, it recognizes that $A$ and $B$ are semantically related, but lacks the document-level temporal context to discern which concept is foundational.

\subsection{Threats to Validity}
\paragraph{Internal Validity:}
The reading-order baseline relies on the first mention of whole-word terms in the text. Concept synonym variation was normalized using canonical aliases, and ties within the same paragraph were treated as unlocated.

\paragraph{Construct Validity:}
AL-CPL Geometry was constructed from Wikipedia concepts and transitively closed by human annotators. In expository textbooks, direct pedagogical dependencies do not always map 1-to-1 to encyclopedic graph links.

\paragraph{External Validity (Pending Experiments):}
Our verified findings are currently established on a single domain (*CK-12 Basic Geometry*). While the underlying reading-order phenomenon is well-supported by educational literature \cite{alzetta2020prereq}, multi-domain replication on *OpenStax Physics* and multi-model scaling (3B vs 14B) remain pending experiments necessary to prove universal cross-domain generalization.

% =========================================================================
\section{Conclusion}
% =========================================================================

This study provides a controlled empirical evaluation of local open-weight LLMs for educational knowledge graph construction. We demonstrate that standalone 7B models fail at directional prerequisite classification, operating near chance ($\text{ROC-AUC} \approx 0.503$). In contrast, a zero-parameter reading-order baseline captures an overwhelming pedagogical signal ($\text{ROC-AUC} = 0.868$), and neuro-symbolic hybrid pipelines derive their directional validity almost entirely from this structural prior. Furthermore, candidate-pair generation is the dominant pipeline bottleneck, discarding over $43\%$ of valid dependencies. Future work must focus on structure-aware candidate generation and multi-domain scaling before local LLMs can autonomously construct curriculum knowledge graphs.

% =========================================================================
% References
% =========================================================================
\begin{thebibliography}{10}

\bibitem[\protect\citeauthoryear{Alzetta et~al.}{2020}]{alzetta2020prereq}
Chiara Alzetta, Alessio Miaschi, and Felice Dell'Orletta.
\newblock 2020.
\newblock Prerequisite relation learning from instructional textbooks: An empirical study across domains and languages.
\newblock In \emph{Proceedings of the 12th Language Resources and Evaluation Conference (LREC)}, pages 4118--4127.

\bibitem[\protect\citeauthoryear{Carta et~al.}{2023}]{carta2023iterative}
Silvio Carta, Alessandro Giuliani, and Gianni Fenu.
\newblock 2023.
\newblock Iterative zero-shot knowledge graph construction with large language models.
\newblock In \emph{Proceedings of the 22nd International Semantic Web Conference (ISWC)}, pages 142--158.

\bibitem[\protect\citeauthoryear{Chen et~al.}{2018}]{chen2018prereq}
Guanliang Chen, Dan Davis, Jie Xie, and Geert-Jan Houben.
\newblock 2018.
\newblock Automated concept prerequisite learning with topological graph constraints.
\newblock \emph{IEEE Transactions on Learning Technologies}, 11(4):456--468.

\bibitem[\protect\citeauthoryear{Ji et~al.}{2021}]{ji2021survey}
Shaoxiong Ji, Shirui Pan, Erik Cambria, Pekka Marttinen, and Philip~S. Yu.
\newblock 2021.
\newblock A survey on knowledge graphs: Representation, acquisition, and applications.
\newblock \emph{IEEE Transactions on Neural Networks and Learning Systems}, 33(2):494--514.

\bibitem[\protect\citeauthoryear{Kommineni et~al.}{2024}]{kommineni2024llm}
Venkata~K. Kommineni, Markus Stocker, and S{\"o}ren Auer.
\newblock 2024.
\newblock Benchmarking open-source large language models for knowledge graph construction.
\newblock \emph{Semantic Web}, 15(3):1--24.

\bibitem[\protect\citeauthoryear{Liang et~al.}{2015}]{liang2015prereq}
Chen Liang, Zhaohui Wu, Wenyi Huang, and C.~Lee Giles.
\newblock 2015.
\newblock Measuring prerequisite relations among concepts.
\newblock In \emph{Proceedings of the 2015 Conference on Empirical Methods in Natural Language Processing (EMNLP)}, pages 1668--1674.

\bibitem[\protect\citeauthoryear{Trajanoska et~al.}{2023}]{trajanoska2023llm}
Marija Trajanoska, Riste Stojanov, and Dimitar Trajanov.
\newblock 2023.
\newblock Evaluating large language models for zero-shot knowledge graph extraction.
\newblock \emph{IEEE Access}, 11:138542--138555.

\end{thebibliography}

\end{document}
